\chapter{Introduction}
\label{chapter:introduction}


\textcolor{red}{The following is my summary—we should go over this together.} \\
Duckietown is an educational and research platform that enables the exploration of autonomous driving concepts using a simple and low-cost environment. 
It also provides Learning Experiences (LX) in which we learned about basic concepts of self-driving cars. 
However, a number of problems arose during the LX that made it difficult or even impossible to control the Duckiebot reliably: \\
Most importantly, our Duckiebot shows \textbf{strong drift}, likely due to hardware issues. 
Because the provided odometry algorithms rely on correct wheel information, we were not able to control our Duckiebot with a PID controller. \\
Also, in the Computer Vision LX, we were supposed to do find lower- and upper-bounds of HSV values by hand to correctly segment lanes in pictures that are taken by the camera of the bot. 
However, we were not able to find good values, that are \textbf{robust under every lighting condtion}. 
As a consequence, the lane following module was not robust. \\
We tried to adress these challenges via \textbf{conditional imitation learning} and making use of a small \textbf{segmentation model}. 
We find that this pipeline enables relyable control of our bot. 
